\chapter{蒙特卡洛仿真与对比实验}

\section{仿真设计}
设置四类实验组：
\begin{itemize}
  \item G1：无噪声+优质初值；
  \item G2：低噪声+中等初值偏差；
  \item G3：中噪声+较差初值；
  \item G4：高噪声+随机初值。
\end{itemize}
每组随机运行 200 次，记录收敛率、平均迭代次数、平均时间与最终误差。

\section{统计指标}
定义成功判据：
\[
\Phi(\theta_{end})<10^{-4},\quad \|\Delta\theta\|_2<10^{-6}.
\]
定义效率指标：
\[
E=\frac{\text{success\_rate}}{\text{mean\_time}}.
\]

\section{实验现象}
从批量运行结果可得：G1、G2 基本稳定收敛；G3 出现部分震荡；G4 在不启用阻尼时失败率明显上升。启用阻尼与自适应步长后，G4 的成功率提升显著，代价是单次运行时间增加。

\section{可视化结论（文字描述）}
\begin{enumerate}
  \item 误差箱线图显示：高斯牛顿分布更集中；
  \item 收敛轨迹图显示：牛顿法前期下降快、后期易震荡；
  \item 参数恢复散点图显示：阻尼策略在高噪声下具有更小方差。
\end{enumerate}

\section{小结}
综合 800 次仿真结果，推荐默认配置为“阻尼高斯牛顿 + 泰勒16阶 + 自适应步长”，在不同噪声等级下均有较稳定表现。
